\subsection{二维随机变量}

\subsubsection{二维离散型随机变量}

\begin{enumerate}
  \item 设二维随机变量$(X,Y)$的分布函数为
  \begin{align*}
    F(x,y)=A\left(B+\arctan\frac{x}{2}\right)\left(C+\arctan\frac{y}{3}\right),
    \quad -\infty<x<+\infty,\ -\infty<y<+\infty
  \end{align*}
  则$A,B,C$分别是$(\quad)$
  \begin{choices}
    \item $\dfrac{1}{\pi^{2}},\ \dfrac{1}{2},\ \dfrac{1}{2}$
    \item $\dfrac{1}{\pi^{2}},\ \dfrac{\pi}{2},\ \dfrac{\pi}{2}$
    \item $\pi^{2},\ \dfrac{1}{2},\ \dfrac{1}{2}$
    \item $\pi^{2},\ \dfrac{\pi}{2},\ \dfrac{\pi}{2}$
  \end{choices}

  \item 已知二维随机变量$(X,Y)$的分布函数为$F(x,y)$,关于$X$与$Y$的边缘分布函数分别是$F_X(x)$及$F_Y(y)$,则概率$P\{X>x_0,\ Y>y_0\}$可表示为$(\quad)$
  \begin{choices}
    \item $F(x_0,y_0)$
    \item $1-F(x_0,y_0)$
    \item $\left[1-F_X(x_0)\right]\cdot\left[1-F_Y(y_0)\right]$
    \item $1-F_X(x_0)-F_Y(y_0)+F(x_0,y_0)$
  \end{choices}

  \item 以下对于二维随机变量$(X,Y)$的分布律与边缘分布律的描述正确的是$(\quad)$
  \begin{choices}
    \item 二维随机变量$(X_1,Y_1)$和$(X_2,Y_2)$的边缘分布相同,则它们必是相同的二维随机变量.
    \item 由二维随机变量$(X,Y)$的分布律必可求得它的边缘分布律.
    \item 由二维随机变量$(X,Y)$的边缘分布律必可求得它的分布律.
    \item 两个边缘分布函数的乘积是联合分布函数.
  \end{choices}

  \item 设二维随机变量$(X,Y)$的分布律如下表,则$P\{XY=0\}=(\quad)$
  \begin{table}[htbp]
    \centering
    \begin{tabular}{cccc}
      \toprule
      $X\backslash Y$ & $0$   & $1$   & $2$   \\
      \midrule
      $0$             & $0.1$ & $0.2$ & $0$   \\
      $1$             & $0.3$ & $0.1$ & $0.1$ \\
      $2$             & $0.1$ & $0$   & $0.1$ \\
      \bottomrule
    \end{tabular}
  \end{table}
  \begin{choices}
    \item $0.8$
    \item $0.7$
    \item $0.1$
    \item $0.3$
  \end{choices}

  \item 以下对于二维随机变量分布函数的描述错误的是$(\quad)$
  \begin{choices}
    \item $F(-\infty,y)=0$
    \item $F(x,-\infty)=0$
    \item $F(+\infty,+\infty)=1$
    \item $F(x,+\infty)=1$
  \end{choices}

  \item 设$X$与$Y$为随机变量,则事件$\{X\leqslant 1,\ Y\leqslant 1\}$的逆事件为
  \begin{choices}
    \item $\{X>1,\ Y>1\}$
    \item $\{X>1,\ Y\leqslant 1\}$
    \item $\{X\leqslant 1,\ Y>1\}$
    \item $\{X>1\}\cup\{Y>1\}$
  \end{choices}

  \item 设$F_1(x)$与$F_2(x)$分别是随机变量$X_1$与$X_2$的分布函数,为使$F(x)=aF_1(x)-bF_2(x)$是某随机变量$X$的分布函数,则应有$(\quad)$
  \begin{choices}
    \item $a=\dfrac{3}{5},\ b=\dfrac{2}{5}$
    \item $a=\dfrac{3}{5},\ b=-\dfrac{2}{5}$
    \item $a=\dfrac{1}{2},\ b=\dfrac{1}{2}$
    \item $a=\dfrac{1}{3},\ b=-\dfrac{1}{3}$
  \end{choices}

  \item 设二维随机变量$(X,Y)$的概率分布如下表,则$a=(\quad)$
  \begin{table}[htbp]
    \centering
    \begin{tabular}{cccc}
      \toprule
      $X\backslash Y$ & $0$      & $1$      & $2$      \\
      \midrule
      $1$             & $1/6$    & $1/4$    & $1/12$   \\
      $2$             & $1/12$   & $a$      & $1/4$    \\
      \bottomrule
    \end{tabular}
  \end{table}
  \begin{choices}
    \item $\dfrac{1}{12}$
    \item $\dfrac{1}{6}$
    \item $\dfrac{1}{4}$
    \item $\dfrac{1}{3}$
  \end{choices}

  \item 从$1,2,3,4$中任取一个数记为$X$,再从$1,\cdots,X$中任取一个数记为$Y$,则$P\{Y=2\}=(\quad)$
  \begin{choices}
    \item $\dfrac{25}{48}$
    \item $\dfrac{13}{48}$
    \item $\dfrac{7}{48}$
    \item $\dfrac{3}{48}$
  \end{choices}

  \item 以下对于二维随机变量分布函数的描述正确的是$(\quad)$
  \begin{choices}
    \item $F(-\infty,-\infty)=0$
    \item $F(-\infty,-\infty)=1$
    \item $F(0,+\infty)=1$
    \item $F(+\infty,-\infty)=1$
  \end{choices}
\end{enumerate}


\subsubsection{二维连续型随机变量}

\begin{enumerate}

 \item \textbf{[多选题]}\ 以下结果哪个正确(\quad)
  \begin{choices}
    \item $\displaystyle\int_{-\infty}^{+\infty}\int_{-\infty}^{+\infty} f(x,y)\,dxdy=1$
    \item $\displaystyle\int_{-\infty}^{+\infty} f_X(x)\,dx=1$
    \item $\displaystyle\int_{-\infty}^{+\infty} f_{Y|X}(y/x)\,dy=1$
    \item $\displaystyle\int_{-\infty}^{+\infty} f_{X|Y}(x/y)\,dx=1$
  \end{choices}

  \item \textbf{[多选题]}\ 设 $(X,Y)$ 为二维离散型随机变量,$P\{X=x_i,\,Y=y_j\}=p_{ij}$,$i,j=1,2,\cdots$.以下结果哪个正确?(\quad)
  \begin{choices}
    \item $\displaystyle\sum_{i=1}^{+\infty}\sum_{j=1}^{+\infty} p_{ij}=1$
    \item $\displaystyle\sum_{i=1}^{+\infty} P\{X=x_i\}=1$
    \item $\displaystyle\sum_{m=1}^{+\infty} P\{Y=y_m\}=1$
    \item $\displaystyle\sum_{i=1}^{+\infty} P(Y=y_i / X=x_k)=1$
  \end{choices}

  \item \textbf{[多选题]}\ 设 $(X,Y)$ 为二维离散型随机变量,$P\{X=x_i,\,Y=y_j\}=p_{ij}$,$i,j=1,2,\cdots$.以下结果哪些正确(\quad)
  \begin{choices}
    \item $\displaystyle P\{X=x_1\}=\sum_{j=1}^{+\infty} p_{1j}$
    \item $\displaystyle P\{X=x_k\}=\sum_{j=1}^{\infty} p_{kj}$
    \item $\displaystyle P(Y=y_2/X=x_1)=\frac{p_{21}}{\sum\limits_{m=1}^{+\infty} p_{1m}}$
    \item $\displaystyle P(Y=y_j/X=x_k)=\frac{p_{kj}}{\sum\limits_{m=1}^{+\infty} p_{km}}$
  \end{choices}

  \item 设 $f(x,y)=\begin{cases} kxy, & 0\le x\le 1,\ 0\le y\le x \\[2pt] 0, & \text{其它} \end{cases}$,则(\quad)
  \begin{choices}
    \item $k=2$
    \item $k=4$
    \item $k=6$
    \item $k=8$
  \end{choices}

  \item 设 $(X,Y)$ 的分布密度函数,$D$ 为平面区域,则(\quad)
  \begin{choices}
    \item $\displaystyle P\big((X,Y)\in D\big)=\iint_D f(x,y)\,dxdy$
    \item $P\big((X,Y)\in D\big)=0$
    \item $P\big((X,Y)\in D\big)=1$
    \item 无法计算
  \end{choices}

  \item 设 $(X,Y)$ 为二维连续型随机变量,$F(x,y)$ 为 $(X,Y)$ 的分布函数,$f(x,y)$ 为 $(X,Y)$ 的分布密度函数,$f_X(x)$ 为关于 $X$ 的边缘密度函数,则对于平面上任一点 $(x_0,y_0)$,有(\quad)
  \begin{choices}
    \item $F(x_0,y_0)=0$
    \item $f(x_0,y_0)=0$
    \item $P(X=x_0,\,Y=y_0)=0$
    \item 都不对
  \end{choices}

  \item 设 $f(x,y)$ 为 $(X,Y)$ 的分布密度函数.以下说法错误的是(\quad)
  \begin{choices}
    \item $f(x,y)$ 为可积函数
    \item $f(x,y)$ 为非负函数
    \item $\displaystyle\int_{-\infty}^{+\infty}\int_{-\infty}^{+\infty} f(x,y)\,dxdy=1$
    \item 对于平面上任一点 $(x,y)$,有 $\dfrac{\partial^2 F}{\partial x\partial y}=f(x,y)$
  \end{choices}

  \item 设 $(X,Y)$ 为二维连续型随机变量,$f(x,y)$ 为 $(X,Y)$ 的分布密度函数.以下结果哪个正确(\quad)
  \begin{choices}
    \item $X$ 是一维连续型随机变量
    \item $X$ 是一维离散型连续型随机变量
    \item 不能判定 $X$ 的类型
    \item 都不对
  \end{choices}

  \item 以下关于边缘分布函数的结果哪个正确(\quad)
  \begin{choices}
    \item $\displaystyle F_{Y|X}(y/x)=\frac{F(x,y)}{f_X(x)}$
    \item $\displaystyle F_{Y|X}(y/x)=\frac{F(x,y)}{f_Y(y)}$
    \item $\displaystyle F_{Y|X}(y/x)=\frac{\int_{-\infty}^{y} f(x,y)\,dy}{\int_{-\infty}^{+\infty} f(x,y)\,dx}$
    \item $\displaystyle F_{Y|X}(y/x)=\frac{\int_{-\infty}^{y} f(x,y)\,dy}{\int_{-\infty}^{+\infty} f(x,y)\,dy}$
  \end{choices}

  \item 以下结果哪个正确(\quad)
  \begin{choices}
    \item $\displaystyle f_X(x)=\int_{-\infty}^{+\infty} f(x,y)\,dy$
    \item $\displaystyle f_X(x)=\int_{-\infty}^{+\infty} f(x,y)\,dx$
    \item $\displaystyle f_X(y)=\int_{-\infty}^{+\infty} f(x,y)\,dx$
    \item 都不对
  \end{choices}

\end{enumerate}

\subsubsection{随机变量的独立性}

\begin{enumerate}
    \item \textbf{[判断题]} 设 $X$ 与 $Y$ 相互独立，则 $2X+3$ 与 $-Y+9X$ 也相互独立 (\quad)
    \begin{choices}
        \item 对
        \item 错
    \end{choices}

    \item \textbf{[判断题]} 设 $X$ 与 $Y$ 相互独立，则 $2X+Y$ 与 $-Y+9$ 也相互独立 (\quad)
    \begin{choices}
        \item 对
        \item 错
    \end{choices}

    \item \textbf{[判断题]} 设 $X$ 与 $Y$ 相互独立，对于任意的区间 $I_1,\ I_2$，事件 $\{X \in I_1\}$ 与 $\{Y \in I_2\}$ 也相互独立。(\quad)
    \begin{choices}
        \item 对
        \item 错
    \end{choices}

    \item \textbf{[判断题]} 设 $(X,Y)$ 为二维连续型随机变量，则 $X$ 与 $Y$ 相互独立的充要条件为 $f(x,y)=f_X(x)f_Y(y)$ (\quad)
    \begin{choices}
        \item 对
        \item 错
    \end{choices}

    \item \textbf{[判断题]} 若 $P(X=x_1,\ Y=y_2)=P(X=x_1)P(Y=y_2)$，则 $X$ 与 $Y$ 相互独立 (\quad)
    \begin{choices}
        \item 对
        \item 错
    \end{choices}

    \item \textbf{[判断题]} 设 $X$ 与 $Y$ 相互独立，则 $2X+3$ 与 $-Y+9$ 也相互独立 (\quad)
    \begin{choices}
        \item 对
        \item 错
    \end{choices}

    \item \textbf{[判断题]} 设 $(X,Y)$ 为二维离散型随机变量，则 $X$ 与 $Y$ 相互独立的充要条件为 $P(X=x_i,\ Y=y_j)=P(X=x_i)P(Y=y_j)$，$i,j=1,2,\cdots$ (\quad)
    \begin{choices}
        \item 对
        \item 错
    \end{choices}

    \item \textbf{[判断题]} 设 $(X,Y,Z)$ 为三维随机变量，若 $X,Y,Z$ 相互独立，则 $X$ 与 $Y$ 也相互独立
    \begin{choices}
        \item 对
        \item 错
    \end{choices}

    \item \textbf{[判断题]} 设 $(X,Y)$ 为二维随机变量，若满足 $F(x,y)=F_X(x)F_Y(y)$，则称 $X$ 与 $Y$ 相互独立 (\quad)
    \begin{choices}
        \item 对
        \item 错
    \end{choices}

    \item \textbf{[判断题]} 设 $X$ 与 $Y$ 相互独立，则 $2X+3Y$ 与 $-Y+9X^2$ 也相互独立 (\quad)
    \begin{choices}
        \item 对
        \item 错
    \end{choices}
\end{enumerate}